\chapter{反抗神明}
在古代地中海地区流传着一个神话。
故事讲述一个名为“闪亮之子”（希伯来语：\texthebrew{הֵילֵל}，希腊语：\textgreek{Φαέθων}）的存在试图篡夺太阳或至高神的位置，结果因其傲慢而坠落或受罚\footnote{%
    译注：在希腊神话中，他试图篡夺太阳神的位置，结果坠落；而在希伯来圣经中，他则试图与至高神同等，结果被罚落入阴间。
}。
这个神话在我们的文化中并不陌生，这主要得益于它的希腊版本。
相比之下，闪米特\footnote{%
    译注：闪米特人是起源于阿拉伯半岛和叙利亚沙漠的游牧民族。
    “闪米特”这一名称源自《圣经》中诺亚的长子 \texthebrew{שֵׁם}（闪）。
}版本则鲜为人知，很大程度上是因为，现存为数不多的书面记载之一在翻译过程中失去了原意：

“明亮之星，早晨之子啊，你何竟从天坠落？你这攻败列国的何竟被砍倒在地上？你心里曾说：‘我要升到天上；我要高举我的宝座在神众星以上；我要坐在聚会的山上，在北方的极处。我要升到高云之上；我要与至上者同等。’然而，你必坠落阴间，到坑中极深之处。”\cite{biblica2011new-international-version}

这里，英文中的“morning star”（晨星\footnote{%
    译注：也即中译本中的“明亮之星”。
}）对应的正是希伯来语 \texthebrew{הֵילֵל}。
由此，我们不难想象这一神话的基本情景：
天空中最明亮的星辰拒绝与其他星星共居一处，在黎明时分、赶在太阳之前跃入天穹。
然而太阳一升起，它便被太阳的光辉抹去。
这种天上的傲慢日复一日，构成一场永无止境的循环。

大家可能更熟悉另一种译法，即《钦定版圣经》\footnote{%
    译注：《钦定版圣经》又称《詹姆斯王圣经》，在苏格兰及英格兰国王詹姆斯六世及一世的命令下翻译的英文版本圣经，于 1611 年出版，自诞生至今一直都是英语国家极受推崇的圣经译本。
}中的“路西法（\textlatin{\emph{Lucifer}}）”，也就是“晨星”在中世纪拉丁语中的名称。

仅此而已。这就是《圣经》中关于路西法的全部故事。
其余都是民间传说，也就是所谓“同人创作”。
例如，路西法曾是天使，他是《约伯记》\footnote{%
    译注：《约伯记》是《希伯来圣经》的第十八卷，《旧约圣经·诗歌智慧书》的第一卷，第一部诗篇性著作。
    约伯这个名字的含义是“仇视的对象”。
}中的撒旦，他是《创世纪》\footnote{%
    译注：《创世纪》是《希伯来圣经》的第一卷书，开篇之作。
    介绍了宇宙的起源（上帝创造天地），人类的起源（上帝创造了亚当和夏娃）和犹太民族的起源，以及犹太民族祖先生活的足迹。
    本书也是上帝全部计划中的开始，展示了上帝的创造是怎样的完美，人类是怎样堕落的，一个民族是如何被上帝拣选发展壮大的。
}中引诱夏娃的蛇，他是《启示录》\footnote{%
    译注：《启示录》是《新约圣经》收录的最后一个作品，写作时间约在公元 90 至 95 年。
    内容主要是对未来的预警，包括对世界末日的预言：接二连三的大灾难，世界朝向毁灭发展的末日光景，并描述最后审判，重点放在耶稣的再来。
}中的兽或龙；
这些内容在圣经中其实并未明确提及。
圣经所呈现的，只是一个骄傲者上升、最后被击落的故事。
（当然，这并不意味着这些后起传统有什么问题。
神圣文本本来就只是宗教这一由观念、行为与制度构成的整体中的一个组成部分\cite{smart1998nature-religion}。）

如果要在《叛逆》中寻找这样一个象征傲慢的角色，那么丘比显然最符合这一位置。
它明确表示，自己之所以将焰安置在结界之中，是为了“干涉”小圆——
换言之，就是试图控制并僭夺这位当下宇宙的创造者。
于是，它最终为自己的罪行付出了代价：遭到了小圆使者焰怒不可遏地惩罚。
影片由此导向一个新的宇宙。
片尾彩蛋似乎暗示着在这个宇宙中，由于焰以一种主动介入的造物主形象存在，丘比的力量几乎已被彻底剥夺。


然而，如果我们考察路西法最著名的“同人创作”之一——
《失乐园》，我们还可以提出另一种解读。
在讨论弥尔顿\footnote{%
    译注：约翰·弥尔顿（1608 年 12 月 9 日—1674 年 11 月 8 日），英国诗人、政论家、民主斗士，被称为英国文学史上伟大的六位诗人之一。
    代表作品有长诗《失乐园》《复乐园》和《力士参孙》。
}的这部长篇史诗之前，首先需要明确“史诗”这一类型的含义。
从支序分类\footnote{%
    译注：术语支序分类借自生物分类学，指按共同祖先及其后代谱系、而非表面相似性来划分类群。
}的角度来看，史诗作为一个类型大多源出自荷马的作品。
通常的定义会从题材和结构两方面概括它的若干类型化特征；
对本文而言，最关键的特征包括：
它书写的是国家、宇宙或世界尺度上的事件；
它围绕一位超凡脱俗、往往拥有超自然力量的英雄展开；
它采用一种区别于日常言说的高昂文体。
此外，史诗通常会以祈请和主题宣告开篇，故事多用拦腰法（\textlatin{\emph{in medias res}}\footnote{%
    译注：\textlatin{\emph{in medias res}} 为拉丁文，意为从事件中间开始，文学上称“拦腰法”。
    这是一种文学与艺术的叙事手法，故事从某个中间点开始展开，而不是从头讲起。
    为了填补背景故事，拦腰法常须辅以倒叙和跳叙，交代早前的事件、纠葛、登场人物与故事背景。
    最常见的情况，就是主线故事的前传、外传或衍生作品。
}），并包含篇幅较长的独白，其中往往至少会有一段倒叙，交代故事开端之前的事件\cite{flannagan1998paradise-lost-introduction}。

《失乐园》之所以格外引人入胜，部分原因在于撒旦这一形象的双重性：
他是一个卑鄙无耻的角色，谎话连篇、欺诈成性、自私自利，并且善于操纵他人；
但与此同时，他又是史诗英雄，其冒险构成了叙事中心。
他极富魅力，也极具说服力，甚至足以让一些重要批评家相信他的事业实际上是正当的——
其中最著名的便是布莱克\footnote{%
    译注：威廉·布莱克（1757 年 11 月 28 日—1827 年 8 月 12 日），英国第一位重要的浪漫主义诗人。
    代表作品有诗集《纯真之歌》《经验之歌》等。
}，他曾认为弥尔顿是“不自知地站在魔鬼一党”\cite{blake1994marriage-heaven-hell}。
然而，撒旦所谓的事业，乃是征服世界，并奴役、灭绝人类！
（这些目标最终全都实现了。
故事结尾，撒旦的儿子兼孙子——“死亡”，以及他的女儿兼妻子——“罪恶”，共同架起了一座从地狱通往人间的桥梁。
人类因“堕落”而受制于“罪恶”，亚当和夏娃也因此终有一死，并将死亡传给自己的后裔，从此人类都注定一死\cite{milton1998paradise-lost}。）

这些矛盾使撒旦在某种意义上成为一个道德暧昧的形象：
从作品结构来看，他是史诗英雄；但在叙事层内，他又无疑是反派。
那么，这一点与《叛逆》有何关系呢？

不妨这样来看，《叛逆》开篇即祈请“圆环之理”，并点明主题：
在一个无可救赎的世界中，绝望不断循环，而遗忘则成为逃离这一循环的方式。
随后，影片将我们直接抛入一个与动画结局完全脱节的情境之中；
观众必须等到影片后段丘比的大段解释，才得知到底发生了什么。
焰作为魔法少女，当然拥有超自然力量，但她的力量并不止于此。
整部影片都发生在她作为魔女所构成的结界之中，因此她又因这一身份而获得更强的力量；
而在此之后，她更因自己的爱而转化为恶魔。
在这一过程中，她也将故事的尺度推向了宇宙层面。
至于那种使其超越一般动画电影通常话语的独特风格，则可参见第十六章。

因此，《叛逆》是一部史诗。
更进一步说，它是一部关于焰如何从第十二集中小圆的“最好的朋友”这一身份，走到影片结尾处称小圆为敌人这一位置的史诗。
换言之，它讲述的是：
在《魔圆》中，小圆是最接近女神的存在，而焰曾是她最亲近、最忠诚的追随者；
在《叛逆》中，焰堕落为恶魔，同时又征服了物质宇宙，并明言要将其引向毁灭。
然而，焰仍然保持着道德上的暧昧性，以至于直到今天，围绕其道德立场的争论仍在各类网络论坛中不绝于耳。

正如动画并不只是讲述一场浮士德式交易，而是在许多方面实际重演了歌德《浮士德》的叙事轨迹，甚至包含了时间旅行这类颇为隐晦的元素；
《叛逆》也不只是一次傲慢的堕落。
无论从结构上看，还是从它将史诗英雄形象与道德堕落相互并置的方式来看，都颇似《失乐园》。
